\begin{theorem}[Ramified Hasse-phase comparison]\label{mod:parameters-ramifiedhassecomparison}
Let $K/F$ be a totally ramified cyclic extension of odd prime degree
$\ell$, with lower break $t$.  Let
\[
 \chi_K=\chi_F\circ N_{K/F},\qquad
 \psi_K=\psi_F\circ\operatorname{Tr}_{K/F},
\]
and suppose
\[
 m_F(\chi_F)=2d+1>1,\qquad 2d+1\geq t+1,
 \qquad m_K(\chi_K)=2d'+1,
\]
with
\[
 2d'+(\ell-1)t=2\ell d,
 \qquad t=0\quad\hbox{or}\quad d\geq t+1.
\]

Fix odd stationary-conductor decompositions and one exact downstairs
\texttt{StationaryClassRepresentative} of the type exposed by the public
\texttt{PhaseReductionResult.representatives} API.  The real stable-high row
maps that supplied representative to the literal same field element upstairs;
inserting these two explicit packages with
\texttt{LocalLamprechtPhaseData.oddOfStationaryClass} gives the actual lower
and upper local phase data.  Their ordered admissible denominators
$\Gamma_F,\Gamma_K$ satisfy the exact conductor identity
\[
 (2d'+1)+n_K(\psi_K)
   =\ell\bigl((2d+1)+n_F(\psi_F)\bigr),
 \qquad \Gamma_K=\Gamma_F\quad\hbox{under }F^\times\longrightarrow K^\times.
\]
The downstairs representative remains a caller choice, not a canonical
stationary representative, and the proved stable-row map makes its upstairs
image agree.  The lower and upper values in
the conclusion are the \texttt{criticalFactor}s of these real local phase
packages, hence the Hasse/refinement phases belonging to the exact choices;
changing either representative must carry its elementary translation at the
same time.  The critical coordinates are tied to the displayed uniformizers
by $\delta_F=\varpi_F^d$ and $\delta_K=\varpi_K^{d'}$.  The supplied residue
equivalence is required pointwise to be the canonical extension residue map.

More precisely, the principal Lean theorem takes the real stable-high datum
\texttt{H : HighStableOddParameterData} and one explicitly supplied
downstairs representative $R$.  Its two sides are definitionally the
\texttt{downstairs} and \texttt{upstairs} fields of
\texttt{H.compatiblePhasePair R}; the upstairs denominator is
\texttt{H.commonDenominator}, and its representative is exactly
\texttt{H.upstairsRepresentative R}.  Thus the public theorem does not ask a
caller to rebuild or identify the real PhaseReduction pair.

Choose uniformizers with
$N_{K/F}(\varpi_K)=\varpi_F$ and explicit residue lifts
$\widetilde x\in\mathcal O_F$.  The comparison certificate records the named
truncation
\[
 N_{K/F}(1+Y_x)
 \equiv 1+\operatorname{Tr}(Y_x)+E_2(Y_x)
       \pmod{\mathfrak p_F^{2d+1}},
 \qquad Y_x=\varpi_K^{d'}\widetilde x,
\]
together with the separate depth statements for $\operatorname{Tr}(Y_x)$
and $E_2(Y_x)$.  It also supplies the unit $\epsilon$ through the exact
specialization
\[
 \operatorname{Tr}_{K/F}(\varpi_K^{2d'})
   =\epsilon\varpi_F^{2d}.
\]
Lean proves from this identity, rather than storing separately, the bilinear
specialization
\[
 \operatorname{Tr}_{K/F}(Y_xY_y)
   =\epsilon\varpi_F^{2d}\widetilde x\widetilde y.
\]
When $t=0$, the finite coordinate $c$ is explicitly the residue, transported
back to $k$, of $\varpi_K^{d'}/\varpi_F^d$.  When $t>0$, the finite coordinate
$\lambda$ is explicitly the transported residue of
\[
 \frac{\sigma(\varpi_K)-\varpi_K}{\varpi_K^{t+1}},
\]
where $\sigma$ is the fixed generator supplied by the prime-cyclic extension
interface.  Thus neither scalar is introduced only through its final
quadratic identity.

For the common residue field, the certificate gives pointwise normal forms
for the actual two critical functions built from the supplied
\texttt{PhaseReduction} choices.  The upper function is transported through
the canonical residue equivalence and, in the tame and characteristic-two
rows, reindexed by the explicitly supplied nonzero scalar $c$.  The
characteristic-two row also retains the separate nonzero coordinate scaling
which identifies its polar character with the normalized absolute-trace
character.  Exactly one of the manuscript's three cases is available: tame odd characteristic with
$t=0$ and $\bar\epsilon=\ell c^2$; wild odd characteristic with $t>0$,
$\ell$ equal to the residue characteristic, and
$\bar\epsilon=-\lambda^{\ell-1}$; or characteristic two with $t=0$ and its
full modulo-eight refinement data.  In the last row the residue cardinality is
recorded as $|k|=2^f$, and the intermediate identity retains the exact factor
\[
 \phi(\ell x)\,\psi_0\!\left(-\binom{\ell}{2}x^2\right)=\phi(x)^\ell
\]
before passing to the modulo-four phase split.  The pointwise identities imply
the phase identities by finite reindexing.  The corresponding finite lemmas
retain every quadratic scalar and affine term and prove
\[
 \boxed{\ \Delta_{3,K}(\chi_K)=\Delta_{3,F}(\chi_F)^\ell\ }.
\]
Here $\Delta_{3,E}$ denotes the phase attached to the supplied denominator
and stationary representative; the statement makes no stronger claim that a
bare Hasse phase is independent of either choice.

The module also isolates three direction-sensitive auxiliary statements.
The Hasse trace lift is from $L$ down to $k$ and has the exact exponent
$[L:k]$ and sign $(-1)^{[L:k]-1}$.  The affine-pencil cancellation is gated
on a real \texttt{ExactOddPhaseAssembly}: its actual conductor is
$m=t+1$ (equivalently the manuscript parameter $b=m-1$ satisfies $b=t$),
$T=t+1$ is odd, and its extension, nonidentity norm-character, identity
twist, and nonidentity twist critical factors are reindexed through an
explicit equivalence.  The extension degree is retained as $[K:F]=p$.
Thus \texttt{AffinePencil} acts only on the exact assembly rows, never on a
free phase tuple.  The
characteristic-two cancellation uses only the positive orientation
$\psi_0(c+c^2)=1$.  A denominator-transport theorem keeps both old and new
ordered denominator pairs, both explicit representative scalings, and the
inverse affine translations.

The comparison is consequently an equality of the real
\texttt{LocalLamprechtPhaseData.criticalFactor}s.  No declaration from the
temporary phase-reduction interface is retained.
\LeanFile{LanglandsFirstMainLemma/Parameters/RamifiedHasseComparison.lean}
\lean{LanglandsFirstMainLemma.ramifiedHasseTraceLiftPhase}
\lean{LanglandsFirstMainLemma.ramifiedHasseTraceLiftPhase_of_oddDegree}
\lean{LanglandsFirstMainLemma.RamifiedHasseOddCriticalData.ofExactAssembly}
\lean{LanglandsFirstMainLemma.RamifiedHasseOddCriticalData.ofExactAssembly_residualQuotient}
\lean{LanglandsFirstMainLemma.ramifiedHasseOddBoundary_affinePencil}
\lean{LanglandsFirstMainLemma.ramifiedHasseArtinSchreier_oriented}
\lean{LanglandsFirstMainLemma.ramifiedHasseCharTwo_binomialCorrection}
\lean{LanglandsFirstMainLemma.RamifiedHasseCharTwoFiniteData.binomialPointwise}
\lean{LanglandsFirstMainLemma.RamifiedHasseCharTwoFiniteData.upperPointwise_of_mod_four_one}
\lean{LanglandsFirstMainLemma.RamifiedHasseCharTwoFiniteData.upperPointwise_of_mod_four_three}
\lean{LanglandsFirstMainLemma.RamifiedHasseComparisonData.traceProductSpecialization}
\lean{LanglandsFirstMainLemma.ramifiedHasseComparison_of_finiteNormalForms}
\lean{LanglandsFirstMainLemma.HighStableOddParameterData.ramifiedHasseComparison_compatiblePhasePair}
\lean{LanglandsFirstMainLemma.ramifiedHasseComparison}
\lean{LanglandsFirstMainLemma.ramifiedHasseComparison_changeDenominator}
\leanok
\SourceText{Theorem thm:ramified-Hasse-comparison.}
\uses{mod:parameters-phasereduction,mod:parameters-stationaryclassundernorm,mod:finitefield-hassefunctionlift,mod:finitefield-affinepencil,mod:finitefield-artinschreier,mod:finitefield-q0powers}
\FormalizationContract The finite normal-form fields state pointwise
identities for the actual Hasse functions constructed from the real local
phase packages, together with the exact case-specific residue and arithmetic
specializations; they do not assume the final power equality.  That equality
is proved by reindexing the actual Hasse sums and applying the finite-field
lemmas.  The principal exported theorem specializes this calculation to the
actual \texttt{HighStableOddParameterData.compatiblePhasePair}; its upstairs
representative and common denominator are constructed by the real stable-high
row.  The real phase-reduction integration uses
\texttt{StationaryClassRepresentative},
\texttt{LocalLamprechtPhaseData.oddOfStationaryClass}, and, for the auxiliary
affine boundary, \texttt{ExactOddPhaseAssembly.residualPhase}.  All quotient
representatives, denominator changes, norm and trace directions, scalar
corrections, and translations remain explicit.
\end{theorem}
